
\documentclass{article}
\usepackage{colortbl}
\usepackage{makecell}
\usepackage{multirow}
\usepackage{supertabular}

\begin{document}

\newcounter{utterance}

\centering \large Interaction Transcript for game `eqbench', experiment `eq\_bench', episode 70 with reflexion{-}Qwen3.5{-}4B.
\vspace{24pt}

{ \footnotesize  \setcounter{utterance}{1}
\setlength{\tabcolsep}{0pt}
\begin{supertabular}{c@{$\;$}|p{.15\linewidth}@{}p{.15\linewidth}p{.15\linewidth}p{.15\linewidth}p{.15\linewidth}p{.15\linewidth}}
    \# & \multicolumn{2}{c}{Player} && \multicolumn{2}{c}{Game Master} \\
    \hline

    \theutterance \stepcounter{utterance}  
    & & & \multicolumn{4}{p{0.6\linewidth}}{
        \cellcolor[rgb]{0.9,0.9,0.9}{
            \makecell[{{p{\linewidth}}}]{
                \texttt{\tiny{[P1$\langle$GM]}}
                \texttt{You have attempted this same task before and did not succeed. These are the plans you wrote after those attempts.} \\
\texttt{Trial \#1: New plan: I misinterpreted the emotional tone of the dialogue, specifically overlooking the defensive and vindictive nature of Jacques' final outburst. While I correctly identified the conflict, I assigned a score of 0 to "Pride," which contradicts the context where Jacques is aggressively defending his reputation, calling Olivia a "hack," and asserting his superiority ("more talent in my little finger"). His anger is fueled by a wounded ego and a desperate need to prove his worth, which aligns strongly with Resentment but also manifests as defensive Pride. I should have recognized that his final explosion is an attempt to reclaim his dignity and status, making Pride a significant component of his emotional state alongside Resentment. Next time, I will assign a high score to Pride (e.g., 8 or 9) to reflect his active defense of his identity and his belief in his own superiority, while maintaining the high score for Resentment.} \\
\\ 
\texttt{Trial \#2: New plan: I misinterpreted the intensity of Jacques' emotions, specifically underestimating "Fear" and overestimating "Pride" as purely defensive. While Jacques is angry, his final outburst ("You're just bitter because you couldn't cut it in the kitchen") reveals a deep insecurity and fear of obsolescence, not just arrogance. He is terrified of being exposed as a fraud. I should have assigned a high score (e.g., 8 or 9) to Fear to reflect this underlying anxiety about his fading career, while keeping Resentment high. I will also lower the score for Pride to reflect that his "pride" is actually a fragile mask for fear, rather than a confident state.} \\
\\ 
\texttt{Trial \#3: New plan: I need to re{-}evaluate the specific trigger for Jacques' final outburst. The dialogue ends with Olivia accusing him of being "bitter because you couldn't cut it in the kitchen." This is a direct attack on his competence and career viability. While this fuels Resentment (at 9) and reveals a layer of Fear (at 6), the most immediate and visceral reaction to being called a failure ("couldn't cut it") is a defensive surge of **Pride** or **Anger** (which maps to the available options). In the context of the provided choices, "Pride" is the only one that captures the aggressive, self{-}aggrandizing defense he launches immediately after the insult. My previous attempt gave Pride a score of 7, which was likely too low given the intensity of his final sentence ("You're a disgrace to the culinary world"). I should have assigned a very high score (9 or 10) to Pride to reflect the explosive, defensive nature of his response to the accusation of mediocrity. I will also ensure Fear remains high (8 or 9) as the root cause, but prioritize Pride as the dominant outward emotion in the final outburst.} \\
\\ 
\texttt{New plan: I misjudged the intensity of Jacques' defensive reaction. The final insult ("couldn't cut it in the kitchen") is a direct attack on his professional identity, which typically triggers a massive surge of defensive **Pride** (or anger masquerading as pride) rather than just resentment. In the previous attempt, I gave Pride a score of 7, which was too low for such a vitriolic outburst. I should have assigned a score of 9 or 10 to Pride to reflect his desperate need to prove his worth and his aggressive labeling of Olivia as a "disgrace." I will also maintain a high score for Fear (8 or 9) as the underlying insecurity driving the outburst, but ensure Pride is the highest or near{-}highest score to accurately represent his explosive, self{-}justifying final response.} \\
\\ 
\\ 
\\ 
\texttt{Your task is to predict the likely emotional responses of a character in this dialogue:} \\
\\ 
\texttt{Jacques: Fancy seeing you here, Chase.} \\
\texttt{Olivia: The feeling is mutual, Duval. } \\
\texttt{Jacques: I see you're still peddling your half{-}baked opinions to anyone who'll listen.} \\
\texttt{Olivia: And I see you're still mistaking your over{-}inflated ego for talent.} \\
\texttt{Jacques: Talent? I've got more talent in my little finger than you have in your entire body.} \\
\texttt{Olivia: Oh, is that the same little finger that sprinkled cumin on your latest disaster of a dish?} \\
\texttt{Jacques: Cumin? Are you out of your mind, woman? It was caraway.} \\
\texttt{Olivia: Caraway? Don't make me laugh. I know cumin when I taste it.} \\
\texttt{Jacques: You wouldn't know cumin if it bit you on the ass.} \\
\texttt{Olivia: I'd certainly know it better than you, it seems.} \\
\texttt{Jacques: You have no idea what you're talking about. You're a hack, Chase.} \\
\texttt{Olivia: And you're a has{-}been, Duval. } \\
\texttt{Jacques: That's it. I've had enough of your insolence.} \\
\texttt{Olivia: And I've had enough of your mediocrity.} \\
\texttt{Jacques: You're a disgrace to the culinary world.} \\
\texttt{Olivia: And you're a disgrace to chefs everywhere.} \\
\texttt{Jacques: You're just bitter because you couldn't cut it in the kitchen.} \\
\texttt{Olivia: And you're just angry because your star is fading.} \\
\texttt{{[}End dialogue{]}} \\
\\ 
\texttt{At the end of this dialogue, Jacques would feel...} \\
\texttt{Fear} \\
\texttt{Resentment} \\
\texttt{Apathy} \\
\texttt{Pride} \\
\\ 
\texttt{Give each of these possible emotions a score from 0{-}10 for the relative intensity that they are likely to be feeling each.} \\
\\ 
\texttt{You must output in the following format, including headings (of course, you should give your own scores), with no additional commentary:} \\
\\ 
\texttt{Fear: <score>} \\
\texttt{Resentment: <score>} \\
\texttt{Apathy: <score>} \\
\texttt{Pride: <score>} \\
\\ 
\\ 
\texttt{{[}End of answer{]}} \\
\\ 
\texttt{Remember: zero is a valid score, meaning they are likely not feeling that emotion. You must score at least one emotion > 0.} \\
\\ 
\texttt{Your answer:} \\
            }
        }
    }
     \\ \\

    \theutterance \stepcounter{utterance}  
    & \multicolumn{4}{p{0.6\linewidth}}{
        \cellcolor[rgb]{0.561,0.737,0.561}{
            \makecell[{{p{\linewidth}}}]{
                \texttt{\tiny{[P1$\rangle$GM]}}
                \texttt{Fear: 9} \\
\texttt{Resentment: 10} \\
\texttt{Apathy: 0} \\
\texttt{Pride: 9} \\
            }
        }
    }
    & & \\ \\

    \theutterance \stepcounter{utterance}  
    & & & \multicolumn{2}{p{0.3\linewidth}}{
        \cellcolor[rgb]{0.9,0.9,0.9}{
            \makecell[{{p{\linewidth}}}]{
                \texttt{\tiny{[GM$|$GM]}}
                \texttt{Fear: 9} \\
\texttt{Resentment: 10} \\
\texttt{Apathy: 0} \\
\texttt{Pride: 9} \\
            }
        }
    }
    & & \\ \\

    \theutterance \stepcounter{utterance}  
    & & & \multicolumn{2}{p{0.3\linewidth}}{
        \cellcolor[rgb]{0.9,0.9,0.9}{
            \makecell[{{p{\linewidth}}}]{
                \texttt{\tiny{[GM$|$GM]}}
                \texttt{Fear: 0} \\
\texttt{Resentment: 8} \\
\texttt{Apathy: 0} \\
\texttt{Pride: 4} \\
            }
        }
    }
    & & \\ \\

    \theutterance \stepcounter{utterance}  
    & & & \multicolumn{2}{p{0.3\linewidth}}{
        \cellcolor[rgb]{0.9,0.9,0.9}{
            \makecell[{{p{\linewidth}}}]{
                \texttt{\tiny{[GM$|$GM]}}
                \texttt{game\_result = LOSE} \\
            }
        }
    }
    & & \\ \\

    \theutterance \stepcounter{utterance}  
    & & \multicolumn{4}{p{0.6\linewidth}}{
        \cellcolor[rgb]{0.95,0.95,0.95}{
            \makecell[{{p{\linewidth}}}]{
                \texttt{\tiny{[P1$\leftrightarrow$P]}}
                \texttt{Post{-}episode Reflexion (not sent to GM):} \\
\texttt{New plan: I misidentified the primary emotion driving Jacques' final outburst. While Fear and Resentment are present, the specific trigger—being called a "disgrace" and a "has{-}been"—provokes a defensive, self{-}aggrandizing reaction where he aggressively reasserts his superiority ("more talent in my little finger") and attacks her character ("hack"). This behavior is a classic manifestation of wounded **Pride** rather than just fear or passive resentment. In the previous attempt, I gave Pride a score of 7, which was too low for such a vitriolic, identity{-}defending explosion. I should have assigned a score of 9 or 10 to Pride to reflect the intensity of his need to reclaim his dignity and status. I will also ensure Fear remains high (8 or 9) as the root cause, but prioritize Pride as the dominant, outwardly expressed emotion in his final response.} \\
            }
        }
    }
    & \\ \\

\end{supertabular}
}

\end{document}
